% Modul Belajar 13, dihasilkan oleh tools/build.py dari tools/konten/p13.py
\documentclass[a4paper]{article}
\usepackage{modulITERA}
\begin{document}
\Sampul{13}{Fungsi dan Modularitas}{Memecah program menjadi bagian bernama: parameter, nilai kembali, prototipe, dan referensi}{CPMK0607 (menerapkan) dan CPMK0608 (menjelaskan) pada konsep dasar algoritma dan sistem komputer}{sekitar 3 jam belajar mandiri, ditambah 150 menit di kelas}{\texttt{02 Hands-on/P13 - Fungsi - Hands-on/}}
\Bagian{Modul 13}{Peta belajar}
Ikuti urutan ini dari atas ke bawah. Setiap bagian memakai apa yang sudah dibahas di bagian sebelumnya.
\par\vspace{4pt}\noindent{\fontsize{9.5pt}{13pt}\selectfont
\begin{tabular}{@{}>{\raggedright\arraybackslash}p{0.140\textwidth}>{\raggedright\arraybackslash}p{0.840\textwidth}@{}}
\kepala{Urutan} & \kepala{Bagian}\\
1 & Tentang modul ini\\\garis
2 & Masalah pembuka: Kode yang disalin berulang\\\garis
3 & Langkah 1: Fungsi void dan pemanggilannya\\\garis
4 & Langkah 2: Parameter dan nilai kembali\\\garis
5 & Langkah 3: Prototipe fungsi\\\garis
6 & Langkah 4: Pass by value dan by reference\\\garis
7 & Langkah 5: Array sebagai parameter\\\garis
8 & Langkah 6: Variabel lokal dan jangkauannya\\\garis
9 & Pesan error yang sering muncul\\\garis
10 & Latihan mandiri\\\garis
11 & Rangkuman\\\garis
12 & Cek pemahaman\\\garis
13 & Exit Ticket dan tugas\\\garis
14 & Bacaan lanjut dan persiapan minggu depan\\
\garisdasar
\end{tabular}}\par\vspace{6pt}
\Bagian{Pembuka}{Tentang modul ini}
Modul ini mendampingi Pertemuan 13. Program kalian makin panjang: membaca data, mengurutkan, mencari, menampilkan. Menaruh semuanya di \texttt{main} membuat program sulit dibaca dan diperbaiki. Minggu ini kalian memecah program menjadi fungsi, bagian kecil bernama yang masing-masing mengerjakan satu tugas.

Isinya sama dengan slide di kelas, tetapi ditulis lebih pelan dan lebih lengkap, supaya kalian bisa mengulang sendiri di rumah tanpa harus menebak maksud slide.

\Sub{Setelah menyelesaikan modul ini, kalian bisa}
\par\vspace{4pt}\noindent{\fontsize{9.5pt}{13pt}\selectfont
\begin{tabular}{@{}>{\raggedright\arraybackslash}p{0.070\textwidth}>{\raggedright\arraybackslash}p{0.680\textwidth}>{\raggedright\arraybackslash}p{0.230\textwidth}@{}}
\kepala{No} & \kepala{Kemampuan} & \kepala{Dibahas di}\\
1 & Menulis program yang dipecah menjadi fungsi dengan parameter dan nilai kembali yang tepat, termasuk fungsi \texttt{void}, fungsi \texttt{bool}, parameter referensi, dan parameter array (Sub-CPMK 13.1, CPMK0607) & Langkah 1 sampai 6\\\garis
2 & Menelusuri aliran nilai pada pemanggilan fungsi dan menjelaskan perbedaan pass by value dan pass by reference serta jangkauan variabel lokal (Sub-CPMK 13.2, CPMK0608) & kotak Tebak dulu, Latihan 3\\
\garisdasar
\end{tabular}}\par\vspace{6pt}
Karena setiap instrumen penilaian dibagi rata antara Bagian A (menulis kode) dan Bagian B (menelusuri dan menjelaskan), yang dinilai bukan hanya program kalian jalan, tetapi juga kalian bisa menjelaskan mengapa program itu jalan.

\Sub{Cara memakai modul ini}
Setiap langkah punya pola yang sama: penjelasan konsep, kadang contoh dari kehidupan sehari-hari, lalu kode yang bisa langsung kalian ketik. Sesudah kode ada kotak \textbf{Tebak dulu}. Tulis tebakan kalian di kertas sebelum membaca \textbf{Jawaban} di bawahnya. Kebiasaan menebak lalu membuktikan inilah yang membuat konsep menempel, bukan sekadar membaca.

\begin{Penting}[Ketik, jangan salin]
Ketik ulang setiap contoh di GDB Online atau editor kalian, jangan disalin-tempel. Saat mengetik, kalian akan membuat salah ketik, membaca pesan error, dan memperbaikinya. Itu bagian dari belajar.
\end{Penting}
\Sub{Yang perlu disiapkan}
Peramban dengan akses ke onlinegdb.com, atau g++ di komputer kalian sendiri. Kertas dan alat tulis untuk menebak keluaran dan membuat trace table. Folder hands-on pertemuan ini dari repo mata kuliah.

\Bagian{Pembuka}{Masalah pembuka: Kode yang disalin berulang}
Program nilai kelas seorang asisten berisi kode mencari maksimum yang disalin empat kali: untuk kuis, tugas, UTS, dan UAS. Suatu hari ia menemukan kesalahan di kode itu.

Sekarang ia harus memperbaiki empat tempat, dan berharap tidak ada yang terlewat.

\begin{MKeluaran}[title={Tampilan yang diinginkan}]
Maks kuis  : 95
Maks tugas : 90
Maks UTS   : 88
Maks UAS   : 92
\end{MKeluaran}
\begin{Tebak}[Coba pikirkan]
Bagaimana caranya menulis kode mencari maksimum cukup sekali, lalu memakainya untuk empat array berbeda?
\end{Tebak}
\begin{Jawaban}[Cara memecahnya]
Temukan yang berulang: Mencari maksimum dilakukan empat kali dengan data berbeda. Beri nama: Jadikan satu subprogram dengan masukan array dan keluaran nilai maksimum. Terjemahkan: Satu fungsi \texttt{int maksimum(int a[], int n)}, dipanggil empat kali.
\end{Jawaban}
\Bagian{Langkah 1}{Fungsi void dan pemanggilannya}
Fungsi adalah sekumpulan perintah yang diberi nama supaya bisa dipanggil berulang kali. Fungsi \texttt{void} mengerjakan sesuatu, misalnya menampilkan, tanpa memberikan nilai kembali.

Saat \texttt{garis(10)} dipanggil, nilai 10 disalin ke parameter \texttt{n}, badan fungsi dijalankan, lalu eksekusi kembali ke baris sesudah pemanggilan.

\begin{Contoh}[Contoh sehari-hari]
Fungsi seperti resep. Resep ditulis sekali di buku, lalu dipakai berkali-kali dengan bahan berbeda. Bahan yang dibawa ke dapur adalah argumen; daftar bahan di resep adalah parameter.
\end{Contoh}
\begin{MKode}{void.cpp}
#include <iostream>
using namespace std;
void garis(int n) {
    for (int i = 0; i < n; i++) cout << "-";
    cout << endl;
}
int main() {
    garis(10);
    cout << "Laporan Nilai" << endl;
    garis(10);
    return 0;
}
\end{MKode}
\begin{Tebak}[]
Sebelum menjalankan program, tulis keluarannya di kertas, karakter demi karakter.
\end{Tebak}
\begin{Jawaban}[]
\begin{MKeluaran}[]
----------
Laporan Nilai
----------
\end{MKeluaran}
\texttt{void} berarti fungsi tidak mengembalikan nilai. \texttt{n} adalah parameter; 10 adalah argumen.
\end{Jawaban}
\Bagian{Langkah 2}{Parameter dan nilai kembali}
Fungsi dengan tipe selain \texttt{void} memberikan nilai kembali lewat \texttt{return}. Pemanggilannya bisa dipakai di mana pun nilai bertipe itu boleh dipakai: disimpan ke variabel, ditampilkan, atau menjadi argumen fungsi lain.

\texttt{return} langsung mengakhiri fungsi. Pastikan setiap jalur di dalam fungsi berakhir dengan \texttt{return}; bila ada jalur yang tidak, g++ memberi peringatan dan hasilnya tidak bisa ditebak.

\begin{MKode}{return.cpp}
double luasLingkaran(double r) {
    return 3.14 * r * r;
}
int lebihBesar(int a, int b) {
    return (a > b) ? a : b;
}
int main() {
    cout << luasLingkaran(10) << endl;
    int m = lebihBesar(7, 12);
    cout << m << " " << lebihBesar(m, 5) << endl;
    return 0;
}
\end{MKode}
\begin{Tebak}[]
Sebelum menjalankan program, tulis keluarannya di kertas, karakter demi karakter.
\end{Tebak}
\begin{Jawaban}[]
\begin{MKeluaran}[]
314
12 12
\end{MKeluaran}
\texttt{return} mengirim nilai ke pemanggil dan langsung mengakhiri fungsi. Setiap jalur fungsi non-void harus berakhir dengan \texttt{return}.
\end{Jawaban}
\Bagian{Langkah 3}{Prototipe fungsi}
Compiler membaca dari atas ke bawah. Fungsi harus dikenalkan sebelum dipanggil. Ada dua cara: tulis definisinya di atas \texttt{main}, atau tulis prototipenya di atas dan definisinya di bawah. Prototipe membuat \texttt{main} tetap di atas dan mudah dibaca.

\begin{MKode}{prototipe.cpp}
#include <iostream>
using namespace std;
bool genap(int x);
int main() {
    for (int i = 1; i <= 4; i++) {
        bool g = genap(i);
        cout << i << (g ? " genap" : " ganjil") << endl;
    }
    return 0;
}
bool genap(int x) {
    return x % 2 == 0;
}
\end{MKode}
\begin{Tebak}[]
Sebelum menjalankan program, tulis keluarannya di kertas, karakter demi karakter.
\end{Tebak}
\begin{Jawaban}[]
\begin{MKeluaran}[]
1 ganjil
2 genap
3 ganjil
4 genap
\end{MKeluaran}
Prototipe memberi tahu compiler nama, tipe kembali, dan parameter fungsi, sehingga definisinya boleh di bawah \texttt{main}.
\end{Jawaban}
\Bagian{Langkah 4}{Pass by value dan by reference}
Secara bawaan, argumen disalin ke parameter (pass by value). Perubahan di dalam fungsi hanya mengenai salinannya. Bila fungsi memang harus mengubah variabel pemanggil, misalnya fungsi tukar, gunakan parameter referensi dengan \texttt{\&}.

\begin{Contoh}[Contoh sehari-hari]
Memberi fotokopi dokumen kepada teman: coretan teman di fotokopi tidak mengubah dokumen asli kalian. Meminjamkan dokumen aslinya berbeda: coretannya langsung tertinggal di dokumen kalian.
\end{Contoh}
\begin{MKode}{referensi.cpp}
void tambahNilai(int x) { x += 10; }
void tambahRef(int &x) { x += 10; }

int main() {
    int a = 5, b = 5;
    tambahNilai(a);
    tambahRef(b);
    cout << a << " " << b << endl;
    return 0;
}
\end{MKode}
\begin{Tebak}[]
Sebelum menjalankan program, tulis keluarannya di kertas, karakter demi karakter.
\end{Tebak}
\begin{Jawaban}[]
\begin{MKeluaran}[]
5 15
\end{MKeluaran}
Tanpa \texttt{\&}, fungsi menerima salinan. Dengan \texttt{\&}, fungsi bekerja langsung pada variabel aslinya.
\end{Jawaban}
\Bagian{Langkah 5}{Array sebagai parameter}
Inilah jawaban masalah pembuka: satu fungsi \texttt{maksimum} dipakai untuk array berbeda. Fungsi tidak tahu panjang array, jadi panjangnya dikirim sebagai parameter kedua.

Berbeda dengan variabel biasa, array tidak disalin. Fungsi bekerja pada array aslinya, sehingga fungsi seperti \texttt{urutkan(a, n)} bisa mengubah isi array pemanggil.

\begin{MKode}{arrayparam.cpp}
int maksimum(int a[], int n) {
    int m = a[0];
    for (int i = 1; i < n; i++)
        if (a[i] > m) m = a[i];
    return m;
}
int main() {
    int kuis[4] = {80, 95, 70, 85};
    int uts[3] = {88, 76, 90};
    cout << maksimum(kuis, 4) << endl;
    cout << maksimum(uts, 3) << endl;
    return 0;
}
\end{MKode}
\begin{Tebak}[]
Sebelum menjalankan program, tulis keluarannya di kertas, karakter demi karakter.
\end{Tebak}
\begin{Jawaban}[]
\begin{MKeluaran}[]
95
90
\end{MKeluaran}
Panjang array dikirim lewat parameter terpisah. Perubahan elemen di dalam fungsi mengenai array aslinya.
\end{Jawaban}
\Bagian{Langkah 6}{Variabel lokal dan jangkauannya}
Setiap variabel hanya bisa dipakai di bagian program tempat ia dideklarasikan (jangkauan atau scope). Variabel lokal sebuah fungsi lahir saat fungsi dipanggil dan hilang saat fungsi selesai. Karena itu dua fungsi boleh memakai nama variabel yang sama tanpa saling mengganggu.

\par\vspace{4pt}\noindent{\fontsize{9.5pt}{13pt}\selectfont
\begin{tabular}{@{}>{\raggedright\arraybackslash}p{0.174\textwidth}>{\raggedright\arraybackslash}p{0.261\textwidth}>{\raggedright\arraybackslash}p{0.261\textwidth}>{\raggedright\arraybackslash}p{0.283\textwidth}@{}}
\kepala{Jenis} & \kepala{Dideklarasikan di} & \kepala{Hidup sampai} & \kepala{Contoh}\\
Parameter & kurung fungsi & fungsi selesai & \texttt{n} pada \texttt{garis(int n)}\\\garis
Lokal fungsi & badan fungsi & fungsi selesai & \texttt{m} pada \texttt{maksimum}\\\garis
Lokal blok & di dalam \texttt{\{ \}} & blok selesai & \texttt{i} pada \texttt{for}\\
\garisdasar
\end{tabular}}\par\vspace{6pt}
\begin{Penting}[]
Variabel \texttt{i} di \texttt{main} dan \texttt{i} di fungsi lain adalah dua variabel berbeda, walaupun namanya sama.
\end{Penting}
\Bagian{Waspada}{Pesan error yang sering muncul}
Semua pesan di tabel ini dihasilkan g++ dengan opsi \texttt{-Wall}, sama seperti yang dipakai \texttt{cek.sh}.
\par\vspace{4pt}\noindent{\fontsize{9.5pt}{13pt}\selectfont
\begin{tabular}{@{}>{\raggedright\arraybackslash}p{0.240\textwidth}>{\raggedright\arraybackslash}p{0.270\textwidth}>{\raggedright\arraybackslash}p{0.470\textwidth}@{}}
\kepala{Kesalahan} & \kepala{Contoh} & \kepala{Pesan pertama dari g++}\\
Fungsi dipanggil sebelum dikenal & \texttt{panggil sebelum definisi} & \texttt{error: 'kuadrat' was not declared in this scope}\\\garis
Tidak ada return di semua jalur & \texttt{int f() tanpa return akhir} & \texttt{warning: control reaches end of non-void function}\\\garis
Banyak argumen tidak cocok & \texttt{lebihBesar(1, 2, 3)} & \texttt{error: too many arguments to function 'int lebihBesar(int, int)'}\\\garis
Literal ke parameter referensi & \texttt{tukar(a, 5)} & \texttt{error: cannot bind non-const lvalue reference of type 'int\&' to an rvalue of ...}\\\garis
Nilai void dipakai & \texttt{int x = garis(5);} & \texttt{error: void value not ignored as it ought to be}\\
\garisdasar
\end{tabular}}\par\vspace{6pt}
Peringatan pada baris kedua patut dianggap serius. Fungsi yang tidak mengembalikan nilai di sebuah jalur memberi hasil acak, dan \texttt{cek.sh} menolaknya karena memakai \texttt{-Werror}.

\Bagian{Latihan}{Latihan mandiri}
Latihan ini sama dengan hands-on di kelas. Semua file ada di \texttt{02 Hands-on/P13 - Fungsi - Hands-on/latihan/}. Kerjakan berurutan, lalu cocokkan dengan \texttt{expected/} atau jalankan \texttt{bash cek.sh}.
\par\vspace{4pt}\noindent{\fontsize{9.5pt}{13pt}\selectfont
\begin{tabular}{@{}>{\raggedright\arraybackslash}p{0.070\textwidth}>{\raggedright\arraybackslash}p{0.360\textwidth}>{\raggedright\arraybackslash}p{0.150\textwidth}>{\raggedright\arraybackslash}p{0.400\textwidth}@{}}
\kepala{No} & \kepala{File} & \kepala{Tingkat} & \kepala{Yang dilatih}\\
1 & \texttt{handson1-bangun.cpp} & Dasar & fungsi \texttt{void} dan fungsi dengan \texttt{return}\\\garis
2 & \texttt{handson2-prima.cpp} & Menengah & fungsi \texttt{bool}, prototipe, array sebagai parameter\\\garis
3 & \texttt{handson3-referensi.cpp} & Tantangan & telusuri parameter, perbaiki tiga kesalahan fungsi\\\garis
+ & \texttt{bonus-statistik.cpp} & Pengayaan & pustaka fungsi untuk array\\
\garisdasar
\end{tabular}}\par\vspace{6pt}
\Sub{Latihan 1: handson1-bangun.cpp}
Tiga fungsi bangun datar dan garis, dipanggil dari \texttt{main}.

\Sub{Latihan 2: handson2-prima.cpp}
Fungsi \texttt{isPrima} dan \texttt{hitungPrima} dengan prototipe di atas \texttt{main}.

\Sub{Latihan 3: handson3-referensi.cpp}
Tiga fungsi dengan kesalahan pass by value, pembagian bulat, dan jalur tanpa \texttt{return}. Telusuri, perbaiki, dan jelaskan.

\Sub{Bonus: bonus-statistik.cpp}
Fungsi minimum, maksimum, rata-rata, dan pengurutan untuk satu array.

\Sub{Latihan menelusuri}
\begin{MKode}{tebak.cpp}
int f(int x, int &y) {
    x = x * 2;
    y = y + x;
    return x + y;
}

int main() {
    int a = 3, b = 4;
    int c = f(a, b);
    cout << a << " " << b << " " << c << endl;
    return 0;
}
\end{MKode}
\begin{Tebak}[]
Tulis keluarannya di kertas tanpa menjalankan program. Buat trace table bila perlu.
\end{Tebak}
\begin{Jawaban}[]
\begin{MKeluaran}[]
3 10 16
\end{MKeluaran}
\texttt{x} salinan dari \texttt{a} menjadi 6, tetapi \texttt{a} tetap 3. \texttt{y} adalah referensi ke \texttt{b}, jadi \texttt{b} menjadi 4 + 6 = 10. Nilai kembali 6 + 10 = 16.
\end{Jawaban}
\Bagian{Penutup}{Rangkuman}
\par\vspace{4pt}\noindent{\fontsize{9.5pt}{13pt}\selectfont
\begin{tabular}{@{}>{\raggedright\arraybackslash}p{0.320\textwidth}>{\raggedright\arraybackslash}p{0.660\textwidth}@{}}
\kepala{Konsep} & \kepala{Yang perlu diingat}\\
Fungsi void dan pemanggilannya & \texttt{void} berarti fungsi tidak mengembalikan nilai.\\\garis
Parameter dan nilai kembali & \texttt{return} mengirim nilai ke pemanggil dan langsung mengakhiri fungsi.\\\garis
Prototipe fungsi & Prototipe memberi tahu compiler nama, tipe kembali, dan parameter fungsi, sehingga definisinya boleh di bawah \texttt{main}.\\\garis
Pass by value dan by reference & Tanpa \texttt{\&}, fungsi menerima salinan.\\\garis
Array sebagai parameter & Panjang array dikirim lewat parameter terpisah.\\\garis
Variabel lokal dan jangkauannya & Variabel \texttt{i} di \texttt{main} dan \texttt{i} di fungsi lain adalah dua variabel berbeda, walaupun namanya sama.\\
\garisdasar
\end{tabular}}\par\vspace{6pt}
\Bagian{Penutup}{Cek pemahaman}
Jawab tanpa menjalankan program, lalu periksa dengan menjalankannya.
\begin{enumerate}
\item Apa beda parameter dan argumen?
\item Mengapa \texttt{void tukar(int x, int y)} tidak menukar variabel pemanggil?
\item Apa gunanya prototipe fungsi?
\item Apakah array di dalam fungsi adalah salinan?
\item Apa yang terjadi pada variabel lokal sesudah fungsi selesai?
\end{enumerate}
\begin{Jawaban}[]
\begin{enumerate}[leftmargin=16pt]
\item Parameter adalah nama di definisi fungsi; argumen adalah nilai yang dikirim saat memanggil.
\item Parameternya salinan; perlu referensi \texttt{int \&x, int \&y}.
\item Mengenalkan fungsi ke compiler sehingga definisinya boleh ditulis sesudah \texttt{main}.
\item Bukan; fungsi bekerja pada array aslinya.
\item Variabel itu hilang dan tidak bisa dipakai lagi.
\end{enumerate}
\end{Jawaban}
\Bagian{Penilaian}{Exit Ticket 13}
Dikerjakan di 25 menit terakhir kelas, sendiri, tanpa AI, internet, atau diskusi. Exit Ticket 13 adalah satu dari 14 Exit Ticket; rata-ratanya menjadi komponen berbobot 25\%.
\par\vspace{4pt}\noindent{\fontsize{9.5pt}{13pt}\selectfont
\begin{tabular}{@{}>{\raggedright\arraybackslash}p{0.080\textwidth}>{\raggedright\arraybackslash}p{0.600\textwidth}>{\raggedright\arraybackslash}p{0.100\textwidth}>{\raggedright\arraybackslash}p{0.200\textwidth}@{}}
\kepala{Soal} & \kepala{Isi} & \kepala{Poin} & \kepala{CPMK}\\
A1 & Tulis fungsi yang mengembalikan nilai terbesar dari tiga bilangan & 25 & CPMK0607\\\garis
A2 & Tulis fungsi \texttt{void} yang mengubah array menjadi kuadrat setiap elemennya & 25 & CPMK0607\\\garis
B1 & Telusuri nilai variabel sebelum dan sesudah pemanggilan fungsi dengan referensi & 20 & CPMK0608\\\garis
B2 & Jelaskan mengapa sebuah fungsi tukar tidak bekerja & 20 & CPMK0608\\\garis
B3 & Jelaskan jangkauan variabel pada potongan kode yang diberikan & 10 & CPMK0608\\
\garisdasar
\end{tabular}}\par\vspace{6pt}
\Bagian{Penilaian}{Tugas 13: Library Matematika}
Kumpulkan sebelum Pertemuan 14 lewat kanal pengumpulan yang diumumkan dosen.

\Sub{Bagian A: program (50 poin)}
Buat program berisi kumpulan fungsi matematika berikut dengan prototipe di atas \texttt{main}, lalu buat menu interaktif untuk memilih dan menguji setiap fungsi dengan validasi masukan.
\begin{enumerate}
\item \texttt{int faktorial(int n)} dan \texttt{int pangkat(int basis, int eksponen)}
\item \texttt{bool isPrima(int n)}
\item \texttt{int gcd(int a, int b)} dan \texttt{int lcm(int a, int b)}
\item \texttt{double luasSegitiga(double alas, double tinggi)} dan \texttt{double kelilingLingkaran(double r)}
\item \texttt{void tampilkanDeret(int n)} yang mencetak 1 sampai n
\item Menu berulang dengan pilihan keluar
\end{enumerate}
\begin{Contoh}[Bonus, tidak wajib]
Tambahkan fungsi \texttt{void tukar(int \&a, int \&b)} dan tunjukkan bedanya dengan versi tanpa \texttt{\&}.
\end{Contoh}
\Sub{Bagian B: penjelasan (50 poin)}
Komentar maksimal 150 kata: telusuri pemanggilan \texttt{lcm(4, 6)} sampai ke \texttt{gcd}, jelaskan kapan memakai \texttt{void} dan kapan memakai tipe kembali, dan beri satu contoh perbedaan pass by value dan pass by reference dari program kalian.

\Sub{Rubrik Tugas 13}
\par\vspace{4pt}\noindent{\fontsize{8.5pt}{11.5pt}\selectfont
\begin{tabular}{@{}>{\raggedright\arraybackslash}p{0.190\textwidth}>{\raggedright\arraybackslash}p{0.070\textwidth}>{\raggedright\arraybackslash}p{0.200\textwidth}>{\raggedright\arraybackslash}p{0.180\textwidth}>{\raggedright\arraybackslash}p{0.170\textwidth}>{\raggedright\arraybackslash}p{0.170\textwidth}@{}}
\kepala{Kriteria} & \kepala{Poin} & \kepala{Sangat baik 85–100} & \kepala{Baik 70–84} & \kepala{Cukup 55–69} & \kepala{Kurang <55}\\
A1 Program berjalan & 20 & tanpa error dan peringatan & ada peringatan & perlu satu perbaikan & tidak terkompilasi\\\garis
A2 Kelengkapan fungsi & 15 & delapan fungsi dan menu & satu fungsi kurang & dua kurang & tiga atau lebih kurang\\\garis
A3 Ketepatan dan validasi & 15 & semua hasil tepat dan masukan divalidasi & satu salah & dua salah & sebagian besar salah\\\garis
B1 Penelusuran pemanggilan & 30 & jejak lcm ke gcd tepat, alasan jelas & satu langkah keliru & tanpa jejak & tidak ada atau disalin\\\garis
B2 Cerita error dan pengujian & 20 & pesan, sebab, perbaikan, uji & pesan tidak dikutip & hanya perbaikan & tidak ada\\
\garisdasar
\end{tabular}}\par\vspace{6pt}
Nilai kriteria = poin × skor level. Contoh: A1 di level Baik dengan skor 80 memberi 20 × 0,80 = 16 poin.

\Bagian{Penutup}{Bacaan lanjut dan persiapan minggu depan}
\begin{enumerate}
\item Deitel, P. dan Deitel, H. C++ How to Program, edisi ke-10, Bab 6.
\item learncpp.com, Bab 2 dan 12 (functions, references).
\item Gaddis, T. Starting Out with C++, Bab 6.
\item Slide \texttt{01 Slide/P13 Fungsi dan Modularitas.pdf} dan hands-on di \texttt{02 Hands-on/P13 - Fungsi - Hands-on/}.
\end{enumerate}
\begin{Penting}[Minggu depan]
Pertemuan 14 membahas rekursi, yaitu fungsi yang memanggil dirinya sendiri, dan struct untuk mengelompokkan data berbeda tipe.
\end{Penting}
\end{document}
